% 解説プリントのコメント（saiten が書く。./tools/kaisetsu.ps1 がこのファイルからプリントを作る）
% 量子力学1 第8章 本編。採点 2026-11-08（スキャン 2026-11-08_08_04_11_PM_からスキャン_p1〜p6。問題6〜8）
%% matome
\begin{daiji}
  \item 位置が確定した状態 $\ket{x'}$ は，運動量の分布 $\abs{\braket{k}{x'}}^{2}=1/(2\pi)$ が一様で，運動量はまったく定まらない。不確定性関係で $\delta x\to0$ なら $\delta p\to\infty$ であり，「運動量が一意に定まる」のではない（問題8 (4)）。
  \item 位置表示では $\matrixel{x}{\hat{p}^{n}}{\psi}=(-\ii\hbar)^{n}\psi^{(n)}(x)$ で，並進演算子 $\hat{T}(a)$ のべき級数は $\psi(x-a)$ のテイラー展開そのものである。「確かめよ」と問われたら，前の小問と別の道筋（ここでは運動量の完全性を使わない道筋）で同じ結果を出す（問題7 (4)）。
\end{daiji}
\begin{maemo}
  \item 運動量の分布を $\abs{\braket{k}{\psi}}^{2}$ で考えない：第7章 問題3 (4)（2026-11-04）で井戸のエネルギーの固有状態について運動量で何がいえるかを答えず，第7章 問題4 (1)（2026-11-07）で $\abs{\tilde{\psi}(k)}^{2}\dd k$ が運動量の確率と読める理由を書かなかった。今回は「エネルギーが一定なら運動量が一意に定まる」と，エネルギーと運動量の確定を結びつけた。「運動量の分布」と問われたら，まず $\braket{k}{\psi}$ を書き，その絶対値の2乗が $k$ によってどう変わるかを見る。
  \item 複素共役とエルミート共役の書き方：第3章 問題2（2026-10-18・2026-10-21）で複素定数の $\abs{A}^{2}$ を $A^{2}$ と書き，第5章 問題7（2026-10-30）で $(\hat{A}\psi)^{*}$ を $\hat{A}^{\dagger}\psi^{*}$ と書いた。今回も（減点はしていないが）問題7 (3) でブラ $\bra{x}\hat{T}(a)$ を $(\hat{T}(a)^{\dagger}\ket{x})^{*}$，$\ket{x-a}^{*}$ と書き，問題7 (5) の位置表示の計算で $(\dd\psi/\dd x)^{*}$ の $*$ を2行落とした（次の行で戻している）。問題6 (5) の $(\partial\tilde{\psi}/\partial k)^{2}$ は $\tilde{\psi}$ が実数なので正しいが，一般には $\abs{\partial\tilde{\psi}/\partial k}^{2}$ である。ケットからブラへは $\dagger$（エルミート共役），数には $*$（複素共役）と使い分け，$\abs{f}^{2}=f^{*}f$ の $*$ を落とさない。
\end{maemo}
%% mondai 7
\begin{kotaeran}
  \machigai{4}{△}{方針} $\matrixel{x}{\hat{p}^{n}}{\psi}=\int(\hbar k)^{n}\braket{x}{k}\braket{k}{\psi}\dd k$ とし，級数を
  \[ \int_{-\infty}^{\infty}\frac{1}{\sqrt{2\pi}}\ee^{\ii k(x-a)}\tilde{\psi}(k)\dd k=\psi(x-a) \]
  にまとめた。計算は正しいが，運動量の完全性を挟んで $\hat{p}^{n}$ を固有値 $(\hbar k)^{n}$ に置きかえたので，(1)〜(3) と同じ道筋になり，問われている「直接位置表示にする」確かめになっていない（テイラー展開の仮定も使っていない）。位置表示の $\matrixel{x}{\hat{p}^{n}}{\psi}=(-\ii\hbar)^{n}\psi^{(n)}(x)$（要点整理）を使えば，第 $n$ 項の係数が $(-\ii a/\hbar)^{n}(-\ii\hbar)^{n}/n!=(-a)^{n}/n!$ になり，級数が $\psi(x-a)$ の $x$ のまわりのテイラー展開と項ごとに一致する。
\end{kotaeran}
\begin{rikai}
  \item 位置表示では $\hat{p}$ は $-\ii\hbar\,\dd/\dd x$ なので，$\hat{T}(a)=\exp(-\ii a\hat{p}/\hbar)$ は $\exp(-a\,\dd/\dd x)$ になる。これは Level 1 の問題1 (3) の $\exp(a\hat{D})f(x)=f(x+a)$ で $a$ を $-a$ にしたものである。運動量表示と位置表示の両方で同じ $\psi(x-a)$ が出ることが「確かめ」の中身である。
  \item (5) で「これでいい？」と書いた交換関係の方法は正しく，解答と同じである。$\hat{T}(a)$ は $\hat{p}$ のべき級数なので $\comm{\hat{p}}{\hat{T}(a)}=0$ で，(3) の $\hat{T}(a)^{\dagger}\hat{T}(a)=\hat{1}$ と合わせれば $\hat{p}$ の期待値は変わらない。
\end{rikai}
%% mondai 8
\begin{kotaeran}
  \machigai{4}{▲}{方針} $\abs{K}^{2}=m/(2\pi\hbar t)$ で $x$ によらないことは正しい。説明に「エネルギーが一定 $\to$ 運動量が一意に定まる。不確定性原理 $\Delta x\,\Delta p\ge\hbar/2$ より，運動量が決まれば位置の確率は全空間に広がる」と書いた。時刻 $0$ の状態 $\ket{x'}$ は位置が確定した状態で，エネルギーの固有状態でも運動量の固有状態でもない。その運動量表示は $\braket{k}{x'}=\ee^{-\ii kx'}/\sqrt{2\pi}$ で，運動量の分布 $\abs{\braket{k}{x'}}^{2}=1/(2\pi)$ はすべての $k$ で同じである。つまり運動量はまったく定まっておらず，不確定性関係の向きが逆である。また，運動量が定まった状態（平面波）は時刻 $0$ からすでに全空間に広がっているので，「局在していた粒子が短い時間の後に広がる」ことの説明にもならない。正しくは，運動量の分布が一様でどれほど大きな運動量の成分も同じ重みで含まれ，大きな運動量の成分ほど速く遠くへ進むので，どれほど短い時間の後にも全空間に広がる，と説明する。
\end{kotaeran}
\begin{rikai}
  \item 位置の固有ケット $\ket{x'}$ の運動量表示 $\braket{k}{x'}$ は絶対値が $k$ によらない。位置の不確定さが $0$ の極限では運動量の不確定さは無限大で，運動量の分布は一様になる（デルタ関数のフーリエ変換が定数になることと同じ）。
  \item 自由粒子では運動量 $\hbar k$ の成分は速さ $\hbar k/m$ で進む（各成分に位相 $\ee^{-\ii\hbar k^{2}t/(2m)}$ が付く）。$\abs{k}$ に上限が無いので，どれほど短い $t>0$ でも，どれほど遠くまで届く成分がある。
  \item 「エネルギーが定まる」と「運動量が定まる」は別のことである。自由粒子のエネルギーの固有状態でも $\pm k$ の重ね合わせがありうる（第7章 問題3）。説明の問題では，まず問題の状態が何の固有状態か（または何の固有状態でもないか）を確かめる。
\end{rikai}
